\scanpage{64}
Les morphismes les plus importants dans la
situation apparaissent dans le diagramme
% The curved leader attaches hexagonal immediately before the live suivant.
\inserted{hexagonal} suivant
\diagram{\begin{tikzcd}[column sep=9mm,row sep=13mm,arrows={line width=.6pt},execute at end picture={
% Upper b: four separate hook components, with unequal lengths.
\draw[line width=.7pt,line cap=round] ([xshift=7.1pt,yshift=6.0pt]s64upper.east) -- ([xshift=8.6pt,yshift=6.6pt]s64upper.east);
\draw[line width=.6pt,line cap=round] ([xshift=5.3pt,yshift=4.6pt]s64upper.east) -- ([xshift=4.6pt,yshift=3.6pt]s64upper.east);
\draw[line width=.6pt,line cap=round] ([xshift=4.3pt,yshift=2.5pt]s64upper.east) .. controls +(0pt,-1.5pt) and +(-1pt,0pt) .. ([xshift=5.4pt,yshift=.8pt]s64upper.east);
\draw[line width=.6pt,line cap=round] ([xshift=6.4pt,yshift=.5pt]s64upper.east) .. controls +(.5pt,.2pt) and +(-.5pt,0pt) .. ([xshift=7.6pt,yshift=.8pt]s64upper.east);
% Central b-prime: four components; the second one is visibly faint.
\draw[line width=.6pt,line cap=round] ([xshift=7.2pt,yshift=6.2pt]s64centre.east) -- ([xshift=8.2pt,yshift=6.8pt]s64centre.east);
\draw[line width=.4pt,black!45,line cap=round] ([xshift=5.9pt,yshift=5.3pt]s64centre.east) -- ([xshift=5.2pt,yshift=4.5pt]s64centre.east);
\draw[line width=.6pt,line cap=round] ([xshift=4.4pt,yshift=3.3pt]s64centre.east) .. controls +(0pt,-1.5pt) and +(-.2pt,.5pt) .. ([xshift=4.4pt,yshift=1.6pt]s64centre.east);
\draw[line width=.6pt,line cap=round] ([xshift=5.3pt,yshift=1.0pt]s64centre.east) -- ([xshift=7.1pt,yshift=.6pt]s64centre.east);
% Upper-right b: three detached pieces at the source.
\draw[line width=.6pt,line cap=round] ([xshift=-8.7pt,yshift=4.8pt]s64right.north west) .. controls +(-.7pt,.9pt) and +(0pt,-.5pt) .. ([xshift=-9.1pt,yshift=6.0pt]s64right.north west);
\draw[line width=.6pt,line cap=round] ([xshift=-8.0pt,yshift=3.8pt]s64right.north west) -- ([xshift=-6.7pt,yshift=3.1pt]s64right.north west);
\draw[line width=.6pt,line cap=round] ([xshift=-5.6pt,yshift=3.0pt]s64right.north west) .. controls +(.6pt,0pt) and +(-.4pt,-.4pt) .. ([xshift=-4.5pt,yshift=3.5pt]s64right.north west);
% Lower b: one detached upper curl, and one lower bend joined to the shaft.
\draw[line width=.7pt,line cap=round] ([xshift=4.7pt,yshift=4.1pt]s64lower.east) .. controls +(0pt,1pt) and +(-.8pt,.4pt) .. ([xshift=6.4pt,yshift=5.3pt]s64lower.east);
\draw[line width=.6pt,line cap=round] ([xshift=4.0pt,yshift=2.4pt]s64lower.east) .. controls +(-.2pt,-2.6pt) and +(-2pt,0pt) .. ([xshift=7pt]s64lower.east);
% s-prime: two unequal visible shafts merge at one head.  Its source curl
% comprises a long U-like part and a separate right arc, with a pen gap.
\draw[->,line width=.85pt,line cap=round] ([xshift=-2.5pt,yshift=8.0pt]s64centre.north west) -- ([yshift=-2pt]s64upper.south east);
\draw[line width=.5pt,line cap=round] ([xshift=.8pt,yshift=9.1pt]s64centre.north west) .. controls +(-3pt,4pt) and +(4pt,-3pt) .. ([yshift=-2pt]s64upper.south east);
\draw[line width=.7pt,line cap=round] ([xshift=-10pt,yshift=7pt]s64centre.north west) .. controls +(0pt,-3pt) and +(-2pt,0pt) .. ([xshift=-3.5pt,yshift=4.5pt]s64centre.north west) .. controls +(1.5pt,0pt) and +(0pt,-1.3pt) .. ([xshift=-2pt,yshift=6.3pt]s64centre.north west);
\draw[line width=.6pt,line cap=round] ([xshift=-.8pt,yshift=5.1pt]s64centre.north west) .. controls +(2.1pt,.6pt) and +(1pt,-1pt) .. ([xshift=1.2pt,yshift=8.0pt]s64centre.north west);
\node[anchor=east] at ([xshift=-18mm,yshift=3pt]s64left.west) {$(40)$};
\node[anchor=west] at ([xshift=34mm,yshift=-3mm]s64bottom.east) {,};
}]
&|[alias=s64upper]|\Sigma^{\unsure{<|*}}(1)\arrow[rr,"\uline{b}",start anchor=east,shorten <=9pt]\arrow[dl,"\uline{s}"']&&|[alias=s64top]|\Sigma&\\
|[alias=s64left]|\Sigma&&|[alias=s64centre]|\Sigma^{*}(2)\arrow[ul,phantom,"\uline{s'}"']\arrow[rr,"\uline{b}'",start anchor=east,shorten <=9pt]\arrow[dl,"\uline{p}"']\arrow[ur,phantom,"\text{cart}" description]\arrow[dr,phantom,"\text{cart}" description]&&|[alias=s64right]|\Sigma^{\unsure{<|*}}(1)\arrow[ul,"\uline{b}"',start anchor=north west,shorten <=7pt]\arrow[dl,"\uline{s}"]\\
&|[alias=s64lower]|\Sigma^{\unsure{<|*}}(1)\arrow[ul,"\uline{s}"']\arrow[rr,"\uline{b}",start anchor=east,shorten <=7pt]&&|[alias=s64bottom]|\Sigma&
\end{tikzcd}}

où les flèches marquées \diagram{\begin{tikzcd}[column sep=large,arrows={line width=.6pt},execute at end picture={
% Four disconnected marks: an upper slant, then three short dot-like pieces.
\draw[line width=.7pt,line cap=round] ([xshift=6.3pt,yshift=6.1pt]s64legend.east) -- ([xshift=8pt,yshift=6.8pt]s64legend.east);
\fill ([xshift=5.1pt,yshift=4.8pt]s64legend.east) circle[radius=.48pt];
\fill ([xshift=5.3pt,yshift=3.1pt]s64legend.east) circle[radius=.48pt];
\fill ([xshift=6pt,yshift=1.7pt]s64legend.east) circle[radius=.43pt];}]
|[alias=s64legend]|{}\arrow[r,start anchor=east,shorten <=8pt]&{}
\end{tikzcd}} (en fait,
$\uline{b},\deleted{\unsure{p_2}},\uline{b'},\uline{s'}$) sont des immersions bordantes \inserted{propres},
% propres is above bordantes; its surrounding curve continues into the
% insertion leader, ending after bordantes before the separate comma.
et les flèches marquées $\longrightarrow$ (en fait,
$\uline{s}$ et $\uline{p}$) sont (dans le cas équisingulier)
fibrantes et propres \unsure{\}} (\unsure{voire|voir}, de plus : fibres
des variétés à bord, dont le bord est
dans le cas de $\unsure{\uline{s}|\uline{s'}}:\Sigma^{\unsure{<|*}}(1)\longrightarrow\Sigma$, l'image
% A short dark raised stroke above this s may be a pen-start or a prime;
% both readings are retained.  No identification with the later s-prime is imposed.
de $\Sigma^{\unsure{<|*}}(2)\overset{\uline{s'}}{-}\Sigma^{*}(1)$
% The initial curve before i.e. continues below the insertion and into its
% long leader.  It is not separately repeated as a literal parenthesis.
\inserted{i.e. l'ens des $(x\prec y)$ tels que $y$ \uline{ne soit pas} successeur de $x$}).
% The insertion is above this map; successeur de x wraps at its right end.
% The following note is oblique in the left margin beside the fibre/boundary passage.
\margin{donc, dans le cas de $\uline{p}:\Sigma^{\unsure{<|*}}(2)\mathrel{\unsure{-|\longrightarrow}}\Sigma^{\unsure{*|<}}(1)$,
le fibré des bords est l'image de $\Sigma^{\unsure{<|*}}(3)$}
On fera attention
que les applications fibrantes $\uline{s},\uline{p}$ n'ont \deleted{\unsure{rien}}\inserted{rien}
de surjectif --- l'image de $\unsure{\uline{s}|s}$ est formé des $x\in\Sigma$
qui \uline{ne sont pas maximaux}, celle de $p$ est formé
des $(x\prec y)$ tels que $y$ ne soit pas maximal.
% Pale bottom-left foliation, another hand; no authorial page number at the head.
\margin{63}
